\section{The negative-test PW limit from an actual adjoint defect}
\label{sec:negative-pw-trace}

\paragraph{MA-AT-002 (the negative-support limit at polynomial weight zero).}
Retain the actual rational spaces and projections of MA-AT-001,
with \(s=0\) and \(0<\alpha<1/2\).  Let \(J\) be the completed
full theta ideal, let \(P_Z=1-P_J\), and put
\[
 d=\tfrac12-\alpha,\qquad u=\tfrac12+\alpha.
 \tag{MA-AT.20}
\]
For a test \(g\in C_c^\infty((-\infty,0))\), define the two causal
tests
\[
 \begin{aligned}
 h(t)&=e^{2\alpha t}\overline{g(-t)},\\
 \ell(t)&=(e^{-2\alpha t}-e^{2\alpha t})
                         \overline{g(-t)} \qquad(t>0),
 \end{aligned}
 \tag{MA-AT.21}
\]
extended by zero elsewhere.  Write
\[
 \begin{aligned}
 \tau_c(g)&=\operatorname{Tr}_H(C_gP_{Z_c}),\\
 \tau_{\rm strip}(g)&=\operatorname{Tr}_H(C_gP_Z),\\
 \tau_+(h)&=\operatorname{Tr}_H(C_hP_{Z_c}).
 \end{aligned}
 \tag{MA-AT.22}
\]
The first trace is the actual ambient trace of MA-AT-001.
The second is the actual MA-WP trace.  The last is independent
of \(c\) by MA-GW-002--003, since \(h\) is causal.
Then
\[
 \boxed{\displaystyle
 \lim_{c\to-\infty}\tau_c(g)
   =\tau_{\rm strip}(g)
        +\overline{\tau_+(h)-\tau_{\rm strip}(h)}.}
 \tag{MA-AT.23}
\]
This equality is first proved by convergence of an actual
trace-class adjoint defect.  Its established spectral evaluations
give the equivalent formula
\[
 \begin{split}
 \lim_{c\to-\infty}\tau_c(g)
   ={}&\sum_{d<\operatorname{Re}\rho<u}
          m_\rho G_0(\rho-\tfrac12)\\
     &+\sum_{\operatorname{Re}\rho\le d}
          m_\rho G_0\bigl(-2\alpha-\overline{\rho-\tfrac12}\bigr),
 \qquad
 G_0(z)=\int g(x)e^{zx}\,dx .
 \end{split}
 \tag{MA-AT.24}
\]
All zeros here and below are nontrivial completed zeros.
The lower boundary is included in the second sum.
Both series converge absolutely.
No conjugation symmetry of a zero divisor is needed in deriving
the reflected exponent.

\begin{proof}
\emph{The exact adjoint kernel.}
The physical weight is \(w(x)=e^{2\alpha|x|}\).
The weighted adjoint of full-space convolution is
\[
 (C_g^*f)(x)
   =\int_0^\infty
       \overline{g(-t)}\,\frac{w(x-t)}{w(x)}f(x-t)\,dt .
 \tag{MA-AT.25}
\]
This follows by changing variables in the weighted inner product,
initially for compact vectors.  The resulting operator is bounded,
as is \(C_g^*\), so the identity extends to \(H\).
Set
\[
                              E=C_g^*-C_h.
 \tag{MA-AT.26}
\]
Choose \(0<a<b\) with \(\operatorname{supp}h\subset[a,b]\).
If \(x\le0\), both \(x\) and \(x-t\) are nonpositive and
\(w(x-t)/w(x)=e^{2\alpha t}\); hence \(Ef(x)=0\).
If \(x>b\), both are positive and the ratio is
\(e^{-2\alpha t}\); hence \(Ef(x)=C_\ell f(x)\).
Thus, as actual operators on \(H\),
\[
 M_{\mathbf1_{(-\infty,0]}}E=0,\qquad
 P_RE=P_RC_\ell\quad(R>b).
 \tag{MA-AT.27}
\]
These are support and kernel identities, not equalities of
scalar trace distributions.

\emph{A fixed compact part and a uniform actual tail.}
For fixed \(R>b\), the operator
\[
                         A_R=M_{\mathbf1_{[0,R]}}E
 \tag{MA-AT.28}
\]
is trace class on the full \(H\).  Its input variable lies in
the compact interval \([-b,R-a]\).
After conjugating by \(\sqrt w\), each of its two terms is a
smooth compact convolution kernel multiplied on the two sides
by bounded compact factors.  This is trace class by the compact
kernel argument of MA-AT-001.  The sharp endpoints and the cusp
of \(w\) cause no difficulty, because those factors are bounded
multipliers of an already trace-class smooth kernel.

The actual projections satisfy
\[
                   P_{Z_c}\longrightarrow P_Z
                   \quad\hbox{strongly as }c\to-\infty.
 \tag{MA-AT.29}
\]
Indeed \(P_c\to1\) strongly, while the increasing actual PW ranges
\(J_c\) exhaust \(J\) after closure, by MA-GP-001--002.
Their orthogonal projections converge strongly to \(P_J\).
This is the actual source condensation in the same physical norm.
Because all these projections are self-adjoint, right
multiplication of the fixed trace-class \(A_R\) gives
\[
             \|A_R(P_{Z_c}-P_Z)\|_1\longrightarrow0.
 \tag{MA-AT.30}
\]

The inclusion \(J_0\subset J\), exactly as in MA-AT.11, also gives
\(P_+P_Z=P_{Z_0}P_+P_Z\), with \(\|P_+P_Z\|\le1\).
Apply the tail factorization to the causal test \(\ell\), both
for \(P_{Z_c}\) and for \(P_Z\).
The one fixed operator \(C_\ell P_{Z_0}\) is trace class by
MA-GW-003.  Therefore, writing
\(\Omega_\ell(R)=\|P_RC_\ell P_{Z_0}\|_1\), one has
\[
 \begin{aligned}
 \|P_REP_{Z_c}\|_1&\le\Omega_\ell(R),\\
 \|P_REP_Z\|_1&\le\Omega_\ell(R),\qquad
                  \Omega_\ell(R)\longrightarrow0.
 \end{aligned}
 \tag{MA-AT.31}
\]
In particular \(EP_{Z_c}\) and \(EP_Z\) are trace class, since
their remaining pieces are \(A_RP_{Z_c}\) and \(A_RP_Z\).
Combining (MA-AT.27)--(MA-AT.31) proves the quantitative estimate
\[
 \|EP_{Z_c}-EP_Z\|_1
   \le\|A_R(P_{Z_c}-P_Z)\|_1+2\Omega_\ell(R)
       \longrightarrow0
 \tag{MA-AT.32}
\]
by first taking \(c\to-\infty\) for fixed \(R\), then
\(R\to\infty\).
Only the defect converges in trace norm here; no such assertion
has been made for \(C_gP_{Z_c}\).

\emph{Adjunction and trace cyclicity after trace class is known.}
MA-AT-001 proves \(C_gP_{Z_c}\in\mathcal S_1\).
Moreover
\(C_g^*P_{Z_c}=C_hP_{Z_c}+EP_{Z_c}\) is trace class by the
causal lift and (MA-AT.31).
The adjoint \(P_{Z_c}C_g^*\) is also trace class.
Cyclicity for these actual products, or their decomposition
into the two diagonal corners of \(P_{Z_c}\), now gives
\[
 \overline{\tau_c(g)}
       =\operatorname{Tr}(C_g^*P_{Z_c})
       =\tau_+(h)+\operatorname{Tr}(EP_{Z_c}).
 \tag{MA-AT.33}
\]
For the actual full-strip projection, MA-WP supplies
\(C_gP_Z,C_hP_Z\in\mathcal S_1\), and (MA-AT.31) supplies
the same defect factor.  Thus the identical justified argument gives
\[
 \overline{\tau_{\rm strip}(g)}
       =\tau_{\rm strip}(h)+\operatorname{Tr}(EP_Z).
 \tag{MA-AT.34}
\]
Subtract (MA-AT.34) from (MA-AT.33), use (MA-AT.32), and
conjugate.  This proves (MA-AT.23) entirely through the actual
operators.

\emph{Evaluation of the already proved trace difference.}
At \(s=0\), MA-GW and MA-WP show that
\(\tau_+(h)-\tau_{\rm strip}(h)\) is the trace over
\(\operatorname{Re}\rho\le d\), with full multiplicities.
For \(\lambda=\rho-\tfrac12\), a direct change of variables gives
\[
 \overline{\int_0^\infty
       e^{2\alpha t}\overline{g(-t)}e^{\lambda t}\,dt}
      =G_0(-2\alpha-\overline{\lambda}).
 \tag{MA-AT.35}
\]
The compact smooth test has rapid vertical Mellin decay,
uniformly on each fixed real strip.  The same is true after
the displayed reflection, so the known unit-height zero count
justifies conjugation and evaluation term by term.
This proves (MA-AT.24).
\end{proof}

\paragraph{The resulting actual joint limit and its local comparison.}
For the source position profiles of MA-AT-001 and every path
\(X,Y\to\infty\), \(Y/X\to0\), \(\delta X\to0\),
the threshold (MA-AT.7) is eventually satisfied.
Since \(g(0)=0\), the actual global trace consequently has limit
\[
 \lim\operatorname{Tr}(M_{\rho_{X,\delta}}C_gQ_Y^{\rm glob})
   =-\tau_{\rm strip}(g)
        -\overline{\tau_+(h)-\tau_{\rm strip}(h)}.
 \tag{MA-AT.36}
\]
This also holds for the fixed annular profile
\(a_X(t)=A(t/X)-A(Xt)\) used in MA-LA-014.
Indeed the output of \(C_gP_{Z_c}\) starts at
\(c+\min\operatorname{supp}g\).  When \(X/Y\) is sufficiently
large, \(A(Xt)\) vanishes on that output, while \(A(t/X)=1\)
through \(x=\log X\).  The same fixed upper-tail estimate applies.
The compact diagonal term for \(P_c\) remains zero.

Choose the fixed finite set \(S\) to contain every prime-power
norm or reciprocal norm met by the test.  The actual local
quotient limit and Meyer's smooth comparison give
\(\operatorname{Tr}T^S_{X,Y,X^{-1}}-P(g)\to-\tau_{\rm all}(g)\).
Subtracting this already established local limit from (MA-AT.36)
gives
\[
 \begin{split}
 &\lim\left\{\operatorname{Tr}(M_{a_X}C_gQ_Y^{\rm glob})
       -\bigl(\operatorname{Tr}T^S_{X,Y,X^{-1}}-P(g)\bigr)\right\}\\
 &\quad=
   \sum_{\operatorname{Re}\rho\ge u}
       m_\rho G_0(\rho-\tfrac12)\\
 &\qquad\quad+
   \sum_{\operatorname{Re}\rho\le d}m_\rho
       \left\{G_0(\rho-\tfrac12)
         -G_0\bigl(-2\alpha-\overline{\rho-\tfrac12}\bigr)\right\}.
 \end{split}
 \tag{MA-AT.37}
\]
At the lower boundary the reflected exponent equals the original
one, so those terms in the second line's lower-divisor correction
cancel exactly.  No exterior zero is asserted to exist.

\paragraph{Scope, characters, and the trace framework.}
The operator proof used only the physical weight, actual nested
supported ranges, their strong condensation, and their proved
causal and full-strip trace-class lifts.
It therefore applies to the actual nontrivial character families
of MA-CGP-001--004 and MA-WP as well, at \(s=0\), with
\(Y=Q\Lambda\), their own divisor, and their actual CQ local
operator; the polar term is then zero.
The conjugate in (MA-AT.35) comes from Hilbert adjunction.
It is not a replacement of a character by its contragredient,
nor does it use a real-coefficient functional equation.

Adding the causal result and this negative result gives an
actual formula for any compact smooth test whose support
avoids a neighbourhood of zero.
This statement does not yet treat every test satisfying only
\(g(0)=0\), and (MA-AT.23)--(MA-AT.37) have not here been
extended to \(s>0\).  At positive polynomial weight the defect
does not vanish on the left, so (MA-AT.32) requires an additional
uniform estimate there.
The finite-length blocks and their lengths remain the spectral
inputs already proved in MA-GW and MA-WP.  The new conclusion is
the actual operator convergence (MA-AT.32), which justifies how
those traces enter the negative-test limit; no categorical trace
class or arithmetic positivity statement replaces that analysis.
